\documentclass[10pt,a4paper]{article}

\usepackage[T1]{fontenc}
\usepackage[utf8]{inputenc}
\usepackage{tgtermes}
\usepackage{tgheros}
\usepackage{tgcursor}
\usepackage{microtype}
\usepackage[a4paper,top=18mm,bottom=20mm,left=18mm,right=18mm]{geometry}
\usepackage{xcolor}
\usepackage{graphicx}
\usepackage{booktabs}
\usepackage{tabularx}
\usepackage{array}
\usepackage{enumitem}
\usepackage[most]{tcolorbox}
\usepackage{fancyhdr}
\usepackage{lastpage}
\usepackage{titlesec}
\usepackage{hyperref}
\usepackage{bookmark}
\usepackage{ragged2e}
\usepackage{caption}
\usepackage{float}
\usepackage{placeins}
\usepackage{needspace}
\usepackage{textcomp}
\usepackage{amsmath}
\usepackage{amssymb}
\usepackage{url}

\definecolor{Navy}{HTML}{17324D}
\definecolor{Teal}{HTML}{008C8C}
\definecolor{TealLight}{HTML}{E1F2F2}
\definecolor{Orange}{HTML}{E58B2A}
\definecolor{OrangeLight}{HTML}{FFF1DE}
\definecolor{Green}{HTML}{2E7D5B}
\definecolor{GreenLight}{HTML}{E8F4ED}
\definecolor{Red}{HTML}{B94A3D}
\definecolor{RedLight}{HTML}{F9E9E7}
\definecolor{Slate}{HTML}{607284}
\definecolor{Light}{HTML}{F2F6F8}
\definecolor{Line}{HTML}{D8E2E9}
\definecolor{Ink}{HTML}{172A3A}

\hypersetup{
  colorlinks=true,
  linkcolor=Teal,
  urlcolor=Teal,
  pdftitle={HeiQ EO vapour cell -- boundary conditions and sizing v2},
  pdfauthor={HeiQ Materials AG | XMembrane project},
  pdfsubject={Quantitative analysis of humidity, temperature, gas films, geometry, operation, and closed-cell behaviour},
  pdfkeywords={XMembrane, electro-osmosis, water vapour, boundary conditions, experiment design, sizing}
}

\graphicspath{{../../outputs/eo_vapor_boundary_analysis_v1/}{../../outputs/eo_vapor_cell_sizing_v1/}}
\setlength{\parindent}{0pt}
\setlength{\parskip}{5pt}
\raggedbottom
\setlist[itemize]{leftmargin=5mm,itemsep=2.5pt,topsep=3pt}
\setlist[enumerate]{leftmargin=6mm,itemsep=3pt,topsep=3pt}
\newcolumntype{Y}{>{\RaggedRight\arraybackslash}X}
\newcolumntype{C}{>{\Centering\arraybackslash}X}
\renewcommand{\arraystretch}{1.20}
\captionsetup{font=small,labelfont={bf,sf,color=Navy}}

\titleformat{\section}
  {\sffamily\bfseries\color{Navy}\fontsize{18}{21}\selectfont}
  {\thesection}{8pt}{}
\titlespacing*{\section}{0pt}{8pt}{6pt}
\titleformat{\subsection}
  {\sffamily\bfseries\color{Navy}\fontsize{12.5}{15}\selectfont}
  {\thesubsection}{7pt}{}

\pagestyle{fancy}
\fancyhf{}
\fancyhead[L]{\sffamily\footnotesize\color{Slate}HEIQ MATERIALS AG · XMEMBRANE}
\fancyhead[R]{\sffamily\footnotesize\color{Slate}EO VAPOUR CELL · V2}
\fancyfoot[L]{\sffamily\scriptsize\color{Slate}Internal technical report · 31 August 2026}
\fancyfoot[R]{\sffamily\scriptsize\color{Slate}Page \thepage\ of \pageref*{LastPage}}
\renewcommand{\headrulewidth}{0.35pt}

\newtcolorbox{mainbox}[1][]{
  enhanced,colback=TealLight,colframe=Teal,boxrule=0.8pt,arc=1.5mm,
  left=4mm,right=4mm,top=3mm,bottom=3mm,#1
}
\newtcolorbox{warningbox}[1][]{
  enhanced,colback=OrangeLight,colframe=Orange,boxrule=0.8pt,arc=1.5mm,
  left=4mm,right=4mm,top=3mm,bottom=3mm,#1
}
\newtcolorbox{stopbox}[1][]{
  enhanced,colback=RedLight,colframe=Red,boxrule=0.8pt,arc=1.5mm,
  left=4mm,right=4mm,top=3mm,bottom=3mm,#1
}
\newtcolorbox{plainbox}[1][]{
  enhanced,colback=Light,colframe=Line,boxrule=0.6pt,arc=1.5mm,
  left=4mm,right=4mm,top=3mm,bottom=3mm,#1
}
\newcommand{\eyebrow}[1]{{\sffamily\bfseries\small\color{Teal}\MakeUppercase{#1}}}
\newcommand{\metric}[3]{%
  \begin{tcolorbox}[colback=white,colframe=Line,boxrule=0.6pt,arc=1.5mm,
  left=3mm,right=3mm,top=3mm,bottom=3mm,height=26mm]
  {\sffamily\bfseries\fontsize{16}{18}\selectfont\color{Navy}#1}\par
  {\sffamily\small\bfseries\color{Ink}#2}\par
  {\sffamily\scriptsize\color{Slate}#3}
  \end{tcolorbox}%
}
\newcommand{\reportfigure}[3]{%
  \begin{figure}[htbp]
    \centering
    \includegraphics[width=#1\textwidth]{#2}
    \caption{#3}
  \end{figure}%
}

\begin{document}

% -----------------------------------------------------------------------------
% Cover
% -----------------------------------------------------------------------------
\thispagestyle{empty}
\pdfbookmark[0]{Cover}{cover}
\vspace*{-10mm}
\begin{tcolorbox}[enhanced,colback=Navy,colframe=Navy,boxrule=0pt,
  arc=0mm,left=11mm,right=11mm,top=13mm,bottom=13mm,width=\textwidth]
  {\sffamily\bfseries\color{TealLight}\large HEIQ XMEMBRANE PROJECT}\par\vspace{7mm}
  {\raggedright\sffamily\bfseries\color{white}\fontsize{29}{33}\selectfont
  Boundary conditions and sizing\par for an EO vapour cell}\par\vspace{4mm}
  {\sffamily\color{white!85}\fontsize{14}{18}\selectfont
  Humidity · temperature · gas films · geometry · drift · closed cell}
\end{tcolorbox}

\vspace{13mm}
{\sffamily\bfseries\fontsize{16}{20}\selectfont\color{Navy}
Internal technical report · v2 · 31 August 2026}\par
\vspace{4mm}
{\fontsize{12}{17}\selectfont\color{Ink}
Quantitative screening of how boundary conditions and the measurement system
affect inference of electro-osmotic water transport in a centimetre-scale
membrane cell.}

\vspace{7mm}
\begin{mainbox}
\textbf{Main conclusion.} Build a modular, once-through fixture around a
50$\times$50 mm membrane aperture, symmetric 2 mm channels, and a fixed
50$\times$50 mm flow field. Use 8--10 mm insulating flow bodies backed by
12 mm metal load plates (or at least 16 mm if the load-bearing plate is PEEK),
and require measured channel-height variation below 20~$\mu$m. The active gas volume is 5.00 mL
per side; the proposed headers, ports, and taps raise the swept internal volume
to 6.22 mL. A small rigid closed cell is a transient demonstrator, not a
steady-state flux measurement.
\end{mainbox}

\vfill
\begin{tabularx}{\textwidth}{@{}p{39mm}Y@{}}
\textbf{Study basis} & 25 cm\textsuperscript{2}, 2 mm gap per side, 32\textdegree C, $a_w=0.50$, 1 m/s, and 6 actual L/min per side \\
\textbf{Fixture basis} & 80$\times$80 mm coupon; 110$\times$110 mm outer plate; 6.22 mL proposed swept volume per side \\
\textbf{Calculations} & 4 material scenarios, 32\,768 combined stress points, and one-factor sweeps \\
\textbf{Numerical control} & 100\% solver success; exact agreement in the registered cross-check points \\
\textbf{Validation level} & Model and experiment design; not material, hardware, or product validation \\
\textbf{Source} & Reproducible Python analysis with configuration, CSV results, figures, and audit file
\end{tabularx}

\clearpage
\pdfbookmark[0]{Executive summary}{summary}
\section*{Executive summary}
\addcontentsline{toc}{section}{Executive summary}

\begin{tcbraster}[raster columns=2,raster equal height=rows,raster column skip=4mm,raster row skip=3mm]
\metric{+10.2\%}{Error from +0.1 RH point}{Single direction, high-diffusivity scenario}
\metric{$-28.9$\%}{Error from +0.1 K}{Single direction, high-diffusivity scenario}
\metric{12.37 mg/h}{Net from 16.80 mg/h gross}{Gas films and back-diffusion leave 73.6\%}
\metric{0.063 s}{Time to 5 Pa}{5 mL rigid cell, gas-only limit}
\end{tcbraster}

\begin{minipage}[t]{0.485\textwidth}
\begin{tcolorbox}[colback=GreenLight,colframe=Green,colbacktitle=Green,
title={\sffamily\bfseries What the analysis shows},fonttitle=\color{white}]
\begin{itemize}
  \item Static humidity and temperature mismatch is highly damaging to single-direction inference.
  \item Stable polarity reversal rejects these static terms almost completely in the model.
  \item Films, gap, flow, and active area are physical transport terms and do not cancel.
  \item ABBA ideally rejects constant and linear time drift.
  \item The closed-cell boundary evolves faster than a typical steady-state criterion.
\end{itemize}
\end{tcolorbox}
\end{minipage}\hfill
\begin{minipage}[t]{0.485\textwidth}
\begin{stopbox}[colbacktitle=Red,title={\sffamily\bfseries What the analysis does not show},fonttitle=\color{white}]
\begin{itemize}
  \item That any material scenario represents the actual membrane.
  \item That the declared ranges are measured probability distributions.
  \item That the 5th/95th stress quantiles are confidence intervals or expected hardware yield.
  \item That ABBA removes hysteresis, curvature, or polarity-linked settling.
  \item That electrode reactions, condensation, or local 2D flow are resolved.
\end{itemize}
\end{stopbox}
\end{minipage}

\begin{warningbox}
\textbf{Decision rule.} Use the once-through cell as the primary measurement
of EO rate. Use the closed cell to demonstrate reversible pumping and to fit a
transient mass balance in which EO, passive back-diffusion, wall/membrane inventory,
and sensor response are estimated together.
\end{warningbox}

\clearpage
\tableofcontents
\clearpage

% -----------------------------------------------------------------------------
\section{Purpose and model boundary}

The purpose is to distinguish three classes of influence:

\begin{enumerate}
  \item \textbf{Physical transport terms:} membrane permeance, EO drag, gas
  films, gap, flow velocity, and active area.
  \item \textbf{Static boundary conditions:} humidity, temperature, and
  pressure difference, edge transport, bypass, and sensor zero offset.
  \item \textbf{Time-dependent measurement terms:} drift, wall/membrane
  inventory, sensor lag, and changes between polarity segments.
\end{enumerate}

For a transparent constant-property limit, the net molar water flux is

\begin{equation}
N_w =
\frac{n_d j/F + K_m(a_{s,b}-a_{r,b})}
{1+K_m\left(k_s^{-1}+k_r^{-1}\right)},
\qquad
K_m=\frac{D_w C_w}{L}.
\label{eq:flux}
\end{equation}

Here $n_dj/F$ is the gross EO flux, $K_m$ is the membrane activity permeance,
and $k_s,k_r$ are the source- and receiver-side film coefficients. The
implemented analysis uses the corresponding nonlinear power-law equation and
iterates the short-channel mean state. The local surrogate reproduces the
project's canonical finite-film solver with a maximum registered relative
difference of $0.0$.

The apparent measured rate is represented schematically by

\begin{equation}
\dot m_{\mathrm{app}} = A_{\mathrm{eff}}M_wN_w
+\dot m_{\mathrm{edge}}+\dot m_{\mathrm{leak}}
+\dot m_{\mathrm{sensor}}+\dot m_{\mathrm{drift}}.
\label{eq:measurement}
\end{equation}

The polarity-odd and polarity-even components are

\begin{equation}
\dot m_{\mathrm{odd}}=\frac{\dot m(+I)-\dot m(-I)}{2},
\qquad
\dot m_{\mathrm{even}}=\frac{\dot m(+I)+\dot m(-I)}{2}.
\label{eq:odd-even}
\end{equation}

Static terms in Equation~\eqref{eq:measurement} primarily enter the
even/single-direction channel. Films, gap, axial depletion, active-area error,
and the residual between polarity segments affect the odd signal.

\subsection{Nominal cell and material scenarios}

\begin{tabularx}{\textwidth}{@{}p{40mm}p{35mm}Y@{}}
\toprule
\textbf{Parameter} & \textbf{Nominal value} & \textbf{Role in the analysis} \\
\midrule
Active area & 25 cm\textsuperscript{2} & 50 $\times$ 50 mm main aperture \\
Channel gap & 2 mm per side & Swept at 1, 1.5, 2, 3, and 4 mm \\
Temperature & 32\textdegree C & Side-to-side mismatch swept from $-0.5$ to $+0.5$ K \\
Water activity & 0.50 & Common nominal inlet; humidity mismatch is swept \\
Pressure & 101\,325 Pa & Differential pressure swept from $-50$ to $+50$ Pa \\
Velocity & 1 m/s & Gives a nominal 6 L/min per side \\
Film correlation & $Sh=7.54$ & Separate 0.25--4 multiplier for unresolved film/support effects \\
\bottomrule
\end{tabularx}

\begin{table}[htbp]
\centering
\caption{Declared material scenarios. These points are not fitted to a named membrane.}
\begin{tabularx}{\textwidth}{@{}Yrrrr@{}}
\toprule
\textbf{Scenario} & $n_d$ & $j$ [A/m\textsuperscript{2}] & $D_w$ [m\textsuperscript{2}/s] & \textbf{Role} \\
\midrule
Low drag, slow $D$ & 0.3 & 10 & $1.60\cdot10^{-14}$ & Signal floor \\
Drag 1, slow $D$ & 1.0 & 10 & $1.60\cdot10^{-14}$ & Order-unity reference \\
Drag 1, fast $D$ & 1.0 & 10 & $1.00\cdot10^{-10}$ & Boundary-sensitive stress point \\
Inverse-design anchor & 100.06 & 1 & $1.60\cdot10^{-14}$ & Hypothetical project signal \\
\bottomrule
\end{tabularx}
\end{table}

% -----------------------------------------------------------------------------
\Needspace{12\baselineskip}
\section{Nominal response and one-factor sensitivity}

\begin{table}[H]
\centering
\caption{Nominal physical response and conditional error in the combined stress design.}
\begin{tabularx}{\textwidth}{@{}Yrrrrr@{}}
\toprule
\textbf{Scenario} & \textbf{Gross} & \textbf{Odd net} & \textbf{Net/gross} & \multicolumn{2}{c}{$q_{0.95}^{\mathrm{stress}}(|e_{\mathrm{odd}}|)$} \\
 & [mg/h] & [mg/h] & [\%] & [mg/h] & [\%] \\
\midrule
Low drag, slow $D$ & 5.041 & 5.041 & 99.99 & 1.25 & 24.80 \\
Drag 1, slow $D$ & 16.804 & 16.803 & 99.99 & 1.25 & 7.47 \\
Drag 1, fast $D$ & 16.804 & 12.369 & 73.61 & 2.10 & 16.98 \\
Inverse-design anchor & 168.146 & 168.136 & 99.99 & 1.71 & 1.01 \\
\bottomrule
\end{tabularx}
\end{table}

The fast-diffusion scenario stands out. High passive permeance makes EO
sensitive to gas films and back-diffusion, so 26.4\% of the gross EO signal is
lost even under nominal boundary conditions. In the slow-diffusion scenarios,
the passive membrane conductance is small and the net value remains close to
the Faradaic term.

\begin{table}[H]
\centering
\caption{Selected one-factor results for the boundary-sensitive fast-diffusion scenario.}
\begin{tabularx}{\textwidth}{@{}Yrr@{}}
\toprule
\textbf{Change from nominal case} & \textbf{Single-direction error} & \textbf{Odd change} \\
\midrule
$+0.1$ RH-percentage-point inlet mismatch & $+10.21$\% & $-0.00001$\% \\
$+0.1$ K side-to-side difference & $-28.87$\% & $-0.00068$\% \\
Gap 2 $\rightarrow$ 1 mm at 1 m/s & $-1.97$\% & $-1.97$\% \\
Gap 2 $\rightarrow$ 4 mm at 1 m/s & $-10.06$\% & $-10.06$\% \\
Flow 6 $\rightarrow$ 3 L/min at 2 mm & $-9.21$\% & $-9.21$\% \\
Flow 6 $\rightarrow$ 12 L/min at 2 mm & $+5.34$\% & $+5.34$\% \\
Film multiplier 1 $\rightarrow$ 0.5 & $-13.98$\% & $-13.98$\% \\
Film multiplier 1 $\rightarrow$ 2 & $+8.85$\% & $+8.85$\% \\
\bottomrule
\end{tabularx}
\end{table}

\reportfigure{0.98}{one_factor_boundary_sensitivity.png}{One-factor sensitivity for the declared high-diffusivity scenario, normalised to its nominal odd rate. Blue circles show apparent single-direction error and orange squares show the polarity-odd result. The humidity and temperature panels apply the same static mismatch in both polarities and therefore show idealised cancellation; they do not demonstrate robustness to mismatch that evolves through an ABBA block. Film, gap, and flow terms remain in the odd estimator. The film multiplier is an unresolved stress coordinate, not a calibrated uncertainty.}
\FloatBarrier

% -----------------------------------------------------------------------------
\section{Humidity, temperature, and polarity sequence}

Two separate humidity generators that both indicate 50\% RH are not
necessarily equal at the membrane. Activity should be calculated from local
vapour pressure and local membrane temperature. A small temperature difference
can therefore appear as a humidity gradient and generate a large passive
contribution in a single-direction measurement.

For an ABBA block, $+I,-I,-I,+I$, the two positive segments and the two
negative segments are averaged in Equation~\eqref{eq:odd-even}. Constant and
linear time drift then have the same average in both polarities and cancel
ideally. The following can remain:

\begin{itemize}
  \item curved or piecewise humidity/temperature drift;
  \item polarity-dependent sensor or electrode settling;
  \item wall and membrane history that depends on the preceding segment;
  \item asymmetric heating or flow changes during polarity reversal;
  \item a timing mismatch between sensor response and the actual water inventory.
\end{itemize}

The global analysis coordinate labelled ``drift between polarities'' must
therefore be interpreted as a \emph{non-cancelling ABBA residual}, not as ideal
linear drift. Its range must be replaced by the residual measured using a foil
blank, a zero-current membrane, and replicated blocks.

At the nominal flow, the model's humidity conversion is approximately
124.5 mg/h per RH percentage point, and the odd estimator sees approximately
half of a residual that arises between the polarity measurements. To keep this
single contribution below the internal 5\% screen, the model gives the
following approximate design targets:

\begin{itemize}
  \item $|\Delta RH_{\mathrm{ABBA,res}}|<0.0041$ percentage points at 5.04 mg/h;
  \item $|\Delta RH_{\mathrm{ABBA,res}}|<0.0135$ percentage points at 16.80 mg/h;
  \item $|\Delta RH_{\mathrm{ABBA,res}}|<0.0099$ percentage points at 12.37 mg/h.
\end{itemize}

These are model-derived control targets, not demonstrated sensor performance.
\FloatBarrier

% -----------------------------------------------------------------------------
\section{Combined stress design and rank association}

The combined design contains 8\,192 Latin-hypercube points per material
scenario, for a total of 32\,768 calculations. It includes small geometric and
electrical errors, humidity/temperature asymmetry, differential pressure,
bypass, edge terms, the film multiplier, and residuals between polarity
segments.

\begin{warningbox}
The ranges are declared engineering stress intervals. Sampling inside the
intervals is chosen for space filling, not because it represents probability
in the future cell. The quantiles and pass fractions below are therefore
conditional on the configuration and are not confidence limits,
probabilities, or expected production yield.
\end{warningbox}

\begin{table}[htbp]
\centering
\caption{Conditional performance in the declared combined stress design. The pass fraction is the fraction of sampled stress points with $|$odd error$|\leq5$\%; it is not a probability or expected hardware yield.}
\begin{tabularx}{\textwidth}{@{}Yrrr@{}}
\toprule
\textbf{Scenario} & $q_{0.95}^{\mathrm{stress}}(|e_{\mathrm{odd}}|)$ & \textbf{Max. $|$odd$|$} & \textbf{Sampled pass fraction} \\
\midrule
Low drag, slow $D$ & 24.80\% & 30.13\% & 20.7\% \\
Drag 1, slow $D$ & 7.47\% & 9.27\% & 67.0\% \\
Drag 1, fast $D$ & 16.98\% & 24.87\% & 41.2\% \\
Inverse-design anchor & 1.01\% & 1.45\% & 100.0\% \\
\bottomrule
\end{tabularx}
\end{table}

\reportfigure{0.91}{combined_sensitivity_rank.png}{Absolute Spearman rank correlation between each input and total odd error, calculated from 8\,192 space-filling stress points per scenario. The values rank monotone associations within the declared ranges; they are not Sobol indices, variance shares, causal effects, or a probability-based uncertainty budget. Small differences between low bars should not be finely ranked without resampling.}

\begin{tabularx}{\textwidth}{@{}p{42mm}Yp{23mm}@{}}
\toprule
\textbf{Scenario} & \textbf{Leading association} & $|\rho|$ \\
\midrule
Low drag, slow $D$ & Residual humidity difference between polarities & 0.98 \\
Drag 1, slow $D$ & Residual humidity difference between polarities & 0.98 \\
Drag 1, fast $D$ & Unresolved film multiplier & 0.73 \\
Inverse-design anchor & Residual humidity difference; then active area & 0.81 / 0.53 \\
\bottomrule
\end{tabularx}
\FloatBarrier

% -----------------------------------------------------------------------------
\section{Films, gap, flow, and area/perimeter}

The project correlation uses $Sh=7.54$ throughout the laminar range. Its local
film coefficient is therefore independent of velocity, while velocity still
changes axial bulk depletion. The real effects of inlet development,
manifolds, support media, and mixing are not known. This is why the analysis
retains an explicit film multiplier.

At 32\textdegree C and 2 mm, the nominal activity-based film coefficient is
$0.09584$ mol m\textsuperscript{$-2$} s\textsuperscript{$-1$} per activity
unit and $Re=247$. In the fast-diffusion scenario, changing the film multiplier
from 0.5 to 2 moves the odd signal from $-14.0$ to $+8.8$\% relative to the
nominal case. Films must therefore be treated as a physical calibration
variable.

The edge-effect model is

\begin{equation}
\dot m_{\mathrm{app}}=J_AA+q_{\mathrm{edge}}P+\dot m_0.
\label{eq:area-perimeter}
\end{equation}

The original 4, 9, and 25 cm\textsuperscript{2} square apertures give three
points for the three parameters in Equation~\eqref{eq:area-perimeter}: zero
residual degrees of freedom and strong area--perimeter correlation.

\begin{table}[htbp]
\centering
\caption{Algebraic design metrics for blank, area, and perimeter terms. The values are not predicted measurement precision.}
\begin{tabularx}{\textwidth}{@{}Yrrrrr@{}}
\toprule
\textbf{Aperture set} & $n$ & \textbf{Residual df} & $\mathrm{corr}(A,P)$ & \textbf{VIF} & \textbf{Condition} \\
\midrule
4/9/25 cm\textsuperscript{2}, square & 3 & 0 & 0.995 & 93.6 & 57.9 \\
+16 cm\textsuperscript{2}, square & 4 & 1 & 0.992 & 62.3 & 54.6 \\
Recommended five-insert design & 5 & 2 & 0.960 & 12.7 & 25.6 \\
\bottomrule
\end{tabularx}
\end{table}

\begin{mainbox}
\textbf{Minimum viable screening design.} Use 20$\times$20, 8$\times$50,
30$\times$30, 40$\times$40, and 50$\times$50 mm apertures. The first two have
the same area (4 cm\textsuperscript{2}) but different perimeters. Run the
8$\times$50 mm slot both parallel and transverse to the flow and replicate the
complete set. Five inserts do not provide robust edge identification by
themselves; the design reduces but does not eliminate shape/perimeter and
hydrodynamic confounding.
\end{mainbox}

% -----------------------------------------------------------------------------
\section{Cell dimensions and mechanical sizing}

\eyebrow{Recommended first build}

The following dimensions are the result of a separate, reproducible sizing
calculation. Geometric and plate-bending relations are physical design
calculations; the expected EO rates still depend on the hypothetical material
scenarios declared earlier in this report.

\begin{table}[H]
\centering
\caption{Recommended baseline dimensions for the first modular fixture.}
\begin{tabularx}{\textwidth}{@{}Yp{37mm}Y@{}}
\toprule
\textbf{Item} & \textbf{Recommended dimension} & \textbf{Quantitative basis} \\
\midrule
Maximum membrane aperture & 50$\times$50 mm (25 cm\textsuperscript{2}) & Largest declared signal and lowest square-mask $P/A=0.80$ cm\textsuperscript{$-1$} \\
Membrane coupon & 80$\times$80 mm & 15 mm land on every side outside the maximum aperture \\
Outer fixture plate & 110$\times$110 mm & A further 15 mm per side for bolts, ports, and load transfer \\
Symmetric gas gap & 2.00 mm per side & 5.00 mL active volume; retain 1 and 4 mm spacers as identification tests \\
Insulating flow body & 8--10 mm & Accommodates ports and leaves at least 4 mm solid behind a 4 mm cavity \\
External metal load plate & 12 mm & Exceeds the plate-screen requirement and provides clamp/port margin \\
All-polymer load plate & at least 16 mm PEEK & Meets the 20~$\mu$m bow screen at the 1 MPa equivalent-load stress point \\
Loaded gap variation & at most 20~$\mu$m & Two per cent of the smallest 1 mm diagnostic gap \\
Seal and guard land & 3 + 2 + 3 mm & Inner seal, vented/tracer guard groove, and outer seal \\
\bottomrule
\end{tabularx}
\end{table}

The 15 mm coupon land accommodates the 8 mm seal/guard stack and leaves 7 mm
before the coupon edge. The second 15 mm land, between the coupon and the outer
plate edge, is reserved for bolts and ports. Hard stops, not gasket compression,
must define the channel height.

For the nominal 50$\times$50$\times$2 mm channel,

\begin{equation}
V_{\mathrm{active}}=Ag
=(50\,\mathrm{mm})(50\,\mathrm{mm})(2\,\mathrm{mm})
=5.00\ \mathrm{mL\ per\ side}.
\end{equation}

The proposed volume budget adds two 7.5 mm-long triangular headers
(0.750 mL), two 6 mm-ID by 8 mm port bores (0.452 mL), and 0.020 mL of
pressure taps. The resulting targets are

\begin{equation}
V_{\mathrm{dead,proposed}}=1.222\ \mathrm{mL},\qquad
V_{\mathrm{cell,swept}}=6.222\ \mathrm{mL\ per\ side}.
\end{equation}

Specify internal dead volume below 1.3 mL per side and measure the assembled
value. Fittings, cell-to-sensor tubing, and instrument volumes are excluded.
At 6 L/min, the proposed internal swept volume corresponds to an ideal 62 ms
residence time, compared with 50 ms for the active parallel-plate region alone.

\subsection{Why the maximum aperture is 25 cm\texorpdfstring{\textsuperscript{2}}{2}}

At fixed current density the gross EO signal scales with area, while an edge
term scales with perimeter:

\begin{equation}
\dot m_{\mathrm{app}}=J_AA+q_PP+\dot m_0,
\qquad
\frac{\dot m_{\mathrm{edge}}}{\dot m_A}
=\frac{q_P}{J_A}\frac{P}{A}.
\end{equation}

The 50 mm square therefore maximises signal and minimises edge sensitivity.
Smaller inserts are identification tools, not alternative nominal cells.

\begin{table}[H]
\centering
\caption{Diagnostic membrane masks. Relative signal assumes equal current density. A 1 mm unrecognised dead edge is shown as a geometric warning, not a prediction.}
\begin{tabularx}{\textwidth}{@{}Yrrrr@{}}
\toprule
\textbf{Exposed membrane shape} & $A$ [cm\textsuperscript{2}] & $P/A$ [cm\textsuperscript{$-1$}] & \textbf{Relative signal} & \textbf{1 mm dead-edge loss} \\
\midrule
20$\times$20 mm square & 4 & 2.00 & 0.16 & 19.0\% \\
8$\times$50 mm slot & 4 & 2.90 & 0.16 & 28.0\% \\
30$\times$30 mm square & 9 & 1.33 & 0.36 & 12.9\% \\
40$\times$40 mm square & 16 & 1.00 & 0.64 & 9.75\% \\
50$\times$50 mm square & 25 & 0.80 & 1.00 & 7.84\% \\
\bottomrule
\end{tabularx}
\end{table}

The recommended inserts mask membrane exposure while retaining the same
50$\times$50 mm gas-flow field, 2 mm gap, and nominal 6 L/min per side. This
keeps the principal channel boundary condition fixed. Run the slot in both
orientations because local edge/support hydrodynamics can still depend on
orientation. If the channel itself is narrowed to follow the mask, the same
slot would require 0.96 or 6.0 L/min at 1 m/s depending on orientation; that
is a deliberately avoided confounding change.

\subsection{Signal mode and film characterisation}

The declared humidity instruments have 0.1 RH-percentage-point expanded
uncertainty per stream. The two-stream expanded uncertainty is therefore
$\sqrt{2}(0.1)=0.141$ RH points, and the declared $S/U=5$ gate requires a
0.707-point differential signal.

\begin{table}[H]
\centering
\caption{Ideal gross-signal detectability for the 25 cm\textsuperscript{2} cell. Passive transport, film loss, storage, leakage, and electrode water are not included.}
\begin{tabularx}{\textwidth}{@{}Yrrrr@{}}
\toprule
\textbf{Reference case} & \textbf{Gross rate} & \textbf{Signal at 6 L/min} & $S/U$ & \textbf{Maximum flow for $S/U=5$} \\
 & [mg/h] & [RH points] & & [L/min per side] \\
\midrule
Low-drag signal floor & 5.041 & 0.0810 & 0.57 & 0.69 \\
Order-unity reference & 16.804 & 0.2699 & 1.91 & 2.29 \\
Inverse-design anchor & 168.146 & 2.7005 & 19.10 & 22.91 \\
\bottomrule
\end{tabularx}
\end{table}

Thus 6 L/min is appropriate as the baseline for film/gap characterisation but
does not resolve the two lower gross signals with the declared humidity
uncertainty. Use a separately calibrated low-flow mode, a higher safe current,
or an integrating donor/receiver water balance. Lowering flow is not a free
sensitivity gain: it also changes films and axial depletion.

\reportfigure{0.98}{gross_signal_envelope_25cm2.png}{Declared gross EO signal for 25 cm\textsuperscript{2} over the configured current and drag coordinates. The figure is an ideal signal envelope, not net flux or material validation.}
\FloatBarrier

\subsection{Wall and load-plate thickness}

For a simply supported square-plate screen under an equivalent uniform load,

\begin{equation}
D_p=\frac{Et^3}{12(1-\nu^2)},\qquad
w_{\max}=\alpha\frac{qa^4}{D_p},\qquad \alpha=0.00406.
\label{eq:plate-screen}
\end{equation}

The span is 50 mm and the allowable bow is 20~$\mu$m. The pressure-only force
on the 25 cm\textsuperscript{2} window is just 0.0125 N at 5 Pa and 0.125 N at
50 Pa, so gas pressure does not determine wall thickness. Bolt and gasket load,
port cut-outs, and the electrode/support stack govern the real fixture.

\begin{table}[H]
\centering
\caption{Required plate thickness for 20~$\mu$m calculated bow. The 200 kPa and 1 MPa rows are equivalent-uniform-load sensitivity coordinates, not calculated gasket pressures.}
\begin{tabularx}{\textwidth}{@{}YrrrY@{}}
\toprule
\textbf{Material} & \textbf{50 Pa gas} & \textbf{200 kPa surrogate} & \textbf{1 MPa surrogate} & \textbf{Selected starting thickness} \\
 & [mm] & [mm] & [mm] & \\
\midrule
PEEK ($E=3.6$ GPa) & 0.57 & 8.98 & 15.35 & 16 mm if fully load-bearing \\
Al 6061-T6 ($E=69$ GPa) & 0.21 & 3.40 & 5.81 & 12 mm external load plate \\
316L ($E=193$ GPa) & 0.15 & 2.43 & 4.16 & 12 mm external load plate \\
\bottomrule
\end{tabularx}
\end{table}

The selected metal thickness is intentionally above the simple plate result to
provide practical margin for fasteners, ports, thermal contact, and imperfect
load distribution. This is still not structural validation. Map the loaded gap
after tightening the real stack, record the torque pattern, and reject a build
that exceeds 20~$\mu$m variation.

\reportfigure{0.98}{cell_dimension_tradeoffs.png}{Left: ideal two-stream humidity signal versus actual flow for the declared gross-rate references; the orange line is the declared $S/U=5$ threshold. Right: plate thickness required to keep calculated bow below 20~$\mu$m across a 50 mm square under equivalent uniform loads. The mechanical load axis is a sensitivity screen, not a gasket contact model.}
\FloatBarrier

% -----------------------------------------------------------------------------
\section{Closed cell: the boundary condition evolves}

In a rigid closed cell, humidity and pressure are not fixed boundary
conditions. With $\delta=a_r-a_s$ and an effective water capacity
$C_{\mathrm{eff}}$ on each side, the linear symmetric limit is

\begin{equation}
\frac{d\delta}{dt}=\frac{2A}{C_{\mathrm{eff}}}
\frac{N_{\mathrm{EO}}-K_m\delta}
{1+K_m(k_s^{-1}+k_r^{-1})}.
\label{eq:closed}
\end{equation}

Walls, seals, the membrane, and sensors can make $C_{\mathrm{eff}}$ much larger
than the gas capacity alone. Until this is measured, the gas-only limit shows
how quickly the pressure constraint can arise.

\begin{table}[htbp]
\centering
\caption{Low-drag scenario at 5.04 mg/h, rigid cell, and gas capacity only.}
\begin{tabularx}{\textwidth}{@{}Yrr@{}}
\toprule
\textbf{Volume per side} & \textbf{Time to 5 Pa} & \textbf{Time to 1 RH point per side} \\
\midrule
5 mL active channel & 0.063 s & 1.21 s \\
6.22 mL proposed swept cell & 0.079 s & 1.50 s \\
1 L & 12.7 s & 4.02 min \\
5 L & 63.4 s & 20.1 min \\
20 L & 254 s & 80.5 min \\
\bottomrule
\end{tabularx}
\end{table}

\reportfigure{0.91}{closed_loop_time_to_5pa.png}{Gas-only time to a 5 Pa differential pressure in an ideal rigid, closed, leak-free cell at $\Delta T=0$ and with no additional wall or membrane capacity. The curves are scenario results, not measured time constants. Sorption, compliance, and pressure balancing can lengthen the response; electrode gas and thermal transients can alter it.}

A symmetric change of one RH percentage point per side corresponds to
approximately 95.2 Pa of water-vapour pressure difference. A temperature
difference of only 0.1 K gives approximately 33.2 Pa of thermal differential
pressure in a rigid cell. This is already far above the $\pm5$ Pa target.

\begin{warningbox}
Do not use a gas leak as passive pressure equalisation: it also transports
water and introduces an unknown bypass term. Use a pressure-balancing bellows,
piston, or active volume control, or accept that the experiment is transient
and fit the complete inventory model.
\end{warningbox}

Leakage, bypass, and pressure-driven transport are not identifiable from the
EO signal alone. They must be measured with a foil blank, tracer tests, and
deliberate pressure challenges. Results obtained at unknown $\Delta P$ must
therefore be reported as a fixture-specific apparent rate, not as intrinsic
membrane flux.
\FloatBarrier

% -----------------------------------------------------------------------------
\Needspace{8\baselineskip}
\section{Recommended experimental sequence}

\eyebrow{Identify the cell before the material}

\begin{enumerate}
  \item \textbf{Foil blank:} run humidity steps, flow, temperature,
  differential pressure, and tracer-leak tests before and after each block.
  Estimate wall memory and sensor lag.
  \item \textbf{Zero-current membrane:} run positive and negative humidity
  gradients to estimate total passive conductance and direction effects.
  \item \textbf{Film/gap:} run 1, 2, and 4 mm gaps and 0.5$\times$,
  1$\times$, and 2$\times$ nominal flow with the same material state.
  \item \textbf{Active ABBA/BAAB:} use one split common inlet, the sequence
  $+I,-I,-I,+I$, and zero-current intervals. Cross-switch the same humidity
  reference between channels.
  \item \textbf{Area/perimeter:} run five aperture types with replicates and
  both orientations of the slot. Fit Equation~\eqref{eq:area-perimeter}.
  \item \textbf{Pressure challenge:} run at least 0, $\pm5$, $\pm20$, and
  $\pm50$ Pa to separate membrane transport from bypass.
  \item \textbf{Closed transient:} perform only after film, passive
  conductance, leakage, and water capacity have been estimated in open-flow
  tests.
\end{enumerate}

For every segment, report source loss, receiver gain, even component, odd
component, current, voltage, local temperature, absolute pressure, and flow
separately. Condensation and Faradaic gases require a separate mass/charge
balance.

% -----------------------------------------------------------------------------
\clearpage
\section{Conclusions and recommended first build}

\begin{mainbox}
\textbf{Build decision.} Use a 50$\times$50 mm maximum membrane aperture in a
fixed 50$\times$50 mm flow field, with 2.00 mm channels on both sides. Use an
80$\times$80 mm membrane coupon, 110$\times$110 mm outer plates, 8--10 mm
insulating flow bodies, and 12 mm metal load plates. If the load-bearing plate
is entirely PEEK, use at least 16 mm as the starting thickness. The proposed
internal volume is 6.22 mL per side, of which 5.00 mL is the active channel and
1.22 mL is headers, ports, and pressure taps.
\end{mainbox}

These dimensions are identical to the canonical recommendation in
\path{outputs/eo_vapor_cell_sizing_v1/audit.json}, which has
\texttt{all\_pass=true}. The geometry is preferred because the
25 cm\textsuperscript{2} aperture gives
the largest available signal and the lowest perimeter-to-area ratio in the
declared centimetre-scale design. At the nominal 2 mm gap and 1 m/s it requires
6 actual L/min per side, has approximately 50 ms active residence time, and has
only 2.9 Pa ideal straight-channel pressure drop. The sizing calculation gives
$Re=236$ using the exact rectangular hydraulic diameter; the boundary model
reports $Re=247$ using the inherited $D_h=2g$ parallel-plate approximation.
This small convention difference is retained explicitly rather than hidden.

The fixture should include:

\begin{itemize}
  \item interchangeable 1, 2, and 4 mm hard-stop spacers;
  \item mask inserts for 20$\times$20, 8$\times$50, 30$\times$30,
  40$\times$40, and 50$\times$50 mm exposed membrane areas while retaining
  the same gas-flow field;
  \item a double seal with a vented or tracer-monitored 2 mm guard groove;
  \item differential-pressure taps at both ends and independent outlet
  back-pressure control;
  \item low-sorption wetted surfaces and symmetric, measured support/electrode
  open area;
  \item an actual-flow range covering at least 3, 6, and 12 L/min per side,
  and up to 24 L/min if the 4 mm/2 m/s stress point is retained.
\end{itemize}

\clearpage
The dimensional recommendations are conditional engineering starting values,
not released manufacturing dimensions. Before material conclusions are drawn,
the assembled fixture must pass the following gates:

\begin{enumerate}
  \item loaded channel-height variation below 20~$\mu$m, documented with the
  real gasket, fastener torque, hard stops, ports, and electrode stack;
  \item measured swept internal dead volume below 1.3 mL per side, with
  external tubing and sensor volumes reported separately;
  \item water-equivalent foil-blank rate below 10\% of the minimum intended EO
  signal, with a 5\% internal target;
  \item main-test cross-membrane pressure within $\pm5$ Pa and quantified
  response at 0, $\pm20$, and $\pm50$ Pa;
  \item donor/receiver mass-balance closure and a measured ABBA/BAAB residual
  that satisfies the signal-specific uncertainty gate;
  \item either less than 5\% change under the gap/flow challenge or a validated
  film correction. Until then, report a geometry-specific effective rate, not
  an intrinsic membrane coefficient.
\end{enumerate}

\clearpage
The nominal 6 L/min humidity-difference measurement does not reach $S/U=5$ for
the 5.04 or 16.80 mg/h gross reference signals. Their ideal flow limits are
0.69 and 2.29 L/min per side, respectively. Use those lower flows only as a
separately calibrated metrology mode; otherwise use integrating water balances
or gravimetry. Finally, closed operation must use external reservoirs with a
bellows, piston, or active volume control and must be fitted as a transient
inventory experiment. Even the proposed 6.22 mL gas-only cell reaches 5 Pa in
about 0.079 s at 5.04 mg/h.

% -----------------------------------------------------------------------------
\section{Reproducibility, audit, and claim boundary}

\begin{tabularx}{\textwidth}{@{}p{50mm}Y@{}}
\toprule
\textbf{Control} & \textbf{Result} \\
\midrule
Combined calculation rows & 32\,768 \\
Solver success fraction & 1.000 \\
Canonical cross-check & Maximum absolute relative difference 0.0 \\
Focused regression and physics tests & 18 passed \\
Sizing audit & \texttt{all\_pass=true}, 10/10 declared checks \\
Calibrated probabilities & No \\
Confidence intervals or expected yield & No \\
Material, hardware, or product validation & No \\
\bottomrule
\end{tabularx}

\begin{stopbox}
All stress quantiles and sampled pass fractions are conditional on the
declared intervals and distribution choices. They are not empirical
probabilities, confidence intervals, or yield forecasts. Gross Faradaic
transport is not net measured transport. The plate loads are sensitivity
coordinates, not predicted gasket pressures. Final acceptance requires
physical blank, mass-balance, flatness, pressure, and leakage tests.
\end{stopbox}

\subsection{Run the analysis}

\begin{verbatim}
cd /home/torsheim/HeiQ/XMembrane_Trond_Study
python3 scripts/analyze_eo_vapor_boundaries_v1.py
python3 scripts/size_eo_vapor_cell_v1.py
python3 -m pytest -q -s \
  tests/test_eo_vapor_boundary_analysis_v1.py \
  tests/test_eo_vapor_cell_sizing_v1.py
\end{verbatim}

\subsection{Compile this report}

\begingroup
\footnotesize
\begin{verbatim}
cd /home/torsheim/HeiQ/XMembrane_Trond_Study
latexmk -cd -pdf -interaction=nonstopmode -halt-on-error \
  reports/eo_vapor_boundary_analysis_v2_en/\
  HeiQ_EO_Vapour_Cell_Boundary_Conditions_and_Sizing_2026-08-31_v2.tex
\end{verbatim}
\endgroup

\subsection{Traceability files}

\begingroup
\footnotesize
\begin{itemize}
  \item Boundary analysis: \path{scripts/analyze_eo_vapor_boundaries_v1.py}
  \item Cell sizing: \path{scripts/size_eo_vapor_cell_v1.py}
  \item Configurations: \path{config/eo_vapor_boundary_analysis_v1.json} and
  \path{config/eo_vapor_cell_sizing_v1.json}
  \item Sizing audit and source hashes:
  \path{outputs/eo_vapor_cell_sizing_v1/audit.json} and
  \path{outputs/eo_vapor_cell_sizing_v1/source_manifest.json}
  \item Boundary-analysis output: \path{outputs/eo_vapor_boundary_analysis_v1/}
  \item Sizing output: \path{outputs/eo_vapor_cell_sizing_v1/}
  \item Tests: \path{tests/test_eo_vapor_boundary_analysis_v1.py} and
  \path{tests/test_eo_vapor_cell_sizing_v1.py}
  \item This report's own freeze:
  \path{reports/eo_vapor_boundary_analysis_v2_en/source_manifest.json}
\end{itemize}
\endgroup

\clearpage
\appendix
\section{Supplementary figures}

\reportfigure{0.98}{activity_temperature_interaction.png}{Interaction between static inlet mismatch and temperature difference for the fast-diffusion scenario. The left panel shows large single-direction bias. The right panel has its own, much smaller scale and shows idealised cancellation when exactly the same static mismatch is used in both polarities; the figure does not include time-dependent ABBA hysteresis.}

\reportfigure{0.98}{combined_error_distributions.png}{Box plots of the space-filling stress samples. These are not empirical or probability-calibrated distributions. The shared scale shows why single-direction inference is vulnerable, but compresses differences in the odd panel; the numerical stress quantiles in the main table are decision-relevant.}
\FloatBarrier

\end{document}
